\documentclass[aps,prd,reprint,superscriptaddress,nofootinbib,amsmath,amssymb,floatfix,longbibliography]{revtex4-2}
\usepackage[T1]{fontenc}
\usepackage{booktabs}
\usepackage{graphicx}
\graphicspath{{figures/}}
\usepackage[hidelinks]{hyperref}

\newcommand{\manuscripttitle}{gravityCU: Multi-GPU VRAM-Resident TreePM Gravity and Full-GPU FoF Solver for Billion-Particle Cosmological \textit{N}-Body Simulations}
\hypersetup{pdftitle={gravityCU: Multi-GPU VRAM-Resident TreePM Gravity and Full-GPU FoF Solver for Billion-Particle Cosmological N-Body Simulations},pdfauthor={Haina Huang}}

\renewcommand{\textfraction}{0.05}
\renewcommand{\topfraction}{0.92}
\renewcommand{\bottomfraction}{0.92}
\renewcommand{\floatpagefraction}{0.85}

\newcommand{\figwide}[3]{%
  \begin{figure*}[!tbp]
    \centering
    \includegraphics[width=168mm]{#1}
    \caption{#2}\label{#3}
  \end{figure*}%
}

\begin{document}
\title{\manuscripttitle}
\author{Haina Huang}
\affiliation{Department of Physics, The Chinese University of Hong Kong, Shatin, Hong Kong, China}
\date{November 26, 2026}

\begin{abstract}
We present \texttt{gravityCU}, a multi-GPU CUDA TreePM, Fast Multipole Method (FMM), and Full-GPU Friends-of-Friends (FoF / Catalog / Spherical Overdensity) solver designed to keep billion-particle ($1024^3 = 1{,}073{,}741{,}824$) cosmological dark-matter simulations fully resident inside device VRAM across $G \in \{1, 2, 3, 4, 6\}$ GPUs with zero host-device particle swapping during time integration. Partitioning the periodic domain into 12 slab-resident or $12\times 12\times 12$ 3D Hilbert/ORB octree zones with cyclic load balancing, dual-stream particle migration, lazy octree refit, zero-spill split-force kernels, and pure-CUDA distributed cuFFT particle-mesh (PM) solvers, four RTX A6000 GPUs retain $139.7$--$142.4\,\mathrm{GB}$ of $1024^3$ state in VRAM while preserving $4.86\,\mathrm{GB}$ minimum free headroom at $z=0$ and $12.31\,\mathrm{GB}$ at $z=1.758$. On dual-socket AMD EPYC PCIe hosts, we resolve a deadlock when direct CUDA peer-to-peer (P2P) access crosses PCIe host bridges (\texttt{NODE} topology) at $1024^3$; gating hardware P2P on NVLink-connected pairs (\texttt{GRAV\_P2P\_MODE=nvlink}) with pinned-host staging delivers a $2.40\times$ multi-GPU speedup at $1024^3$ and $45.4\times$ ($17.53\,\mathrm{s}$ vs.\ $796.36\,\mathrm{s}$, or $50.6\times$ excluding export I/O) on $256^3$ 3D Hilbert execution relative to a 256-thread CPU \texttt{SWIFT} baseline, while GPU FoF catalog ($12.48\times$) and spherical-overdensity ($15.33\times$) post-processing yield zero halo-ID mismatches out of $16.78\mathrm{M}$ particles. At $1024^3$, four RTX A6000 GPUs achieve $7.26\times$ evolution speedup at $z = 0.063 \rightarrow 0$ ($4{,}388.62\,\mathrm{s}$ vs.\ $31{,}878.77\,\mathrm{s}$) and $9.27\times$ solo ($7.46\times$ shared node) at $z = 2.033 \rightarrow 1.758$ ($3{,}510.98\,\mathrm{s}$ vs.\ $32{,}538.37\,\mathrm{s}$), matching $P_\delta(k)$ to $0.00453\%$--$0.01019\%$ and synchronized $P_\theta(k)$ to $0.0111\%$--$0.0405\%$. Reference-independent analytical and conservation gates pass across a 1D Zel'dovich plane wave ($\Delta x_{\rm RMS}/L = 4.4\times 10^{-4}$, $\Delta v_{\rm RMS}/v_{\rm peak} = 0.59\%$), exact periodic Ewald $\theta$-convergence ($1.14\times 10^{-4}$ short-force median at $\theta=0.3$), 32K isolated Plummer virial stability ($|\Delta E/E_0| = 1.15\times 10^{-5}$), and full $256^3$ $z = 200 \rightarrow 0$ Layzer--Irvine closure ($0.0307\%$). Finally, a controlled same-node benchmark against \texttt{CPU SWIFT} ($256\mathrm{T}$), \texttt{GADGET-4} ($128$ cores), and \texttt{PKDGRAV3} ($1\times$/$2\times$ A6000) using $200\,\mathrm{ms}$ NVML hardware telemetry demonstrates up to $49.62\times 10^6$ active-particle updates/s, $20.05\,\mathrm{J}$ GPU ($23.64\,\mathrm{J}$ node) per million updates, and $22.0$--$25.8\times$ whole-node energy reduction.
\end{abstract}

\maketitle

\section{Introduction}\label{sec:intro}

Cosmological $N$-body simulations remain the foundational numerical laboratory for nonlinear structure formation, halo mass functions, and precision summary statistics required by Stage-IV dark-energy and large-scale-structure surveys. Modern hybrid gravity solvers combine a fast Fourier particle-mesh (PM) long-range force on coarse grids~\citep{Hockney1988} with hierarchical spatial decompositions---either Barnes--Hut multipole trees or Cartesian Fast Multipole Methods (FMM)~\citep{Dehnen2000FMM,Dehnen2014FMM}---as implemented in flagship CPU and GPU-accelerated production codes including \texttt{GADGET-2}~\citep{Springel2005GADGET2}, \texttt{GADGET-4}~\citep{Springel2021GADGET4}, \texttt{PKDGRAV3}~\citep{Potter2017PKDGRAV3}, \texttt{HACC}~\citep{Habib2016HACC}, \texttt{ABACUS}~\citep{Garrison2021ABACUS}, and task-based \texttt{SWIFT}~\citep{Schaller2024SWIFT}.

While GPU kernels dramatically accelerate short-range P2P/M2L/M2P force evaluations, scaling to $N = 1024^3 \approx 1.074 \times 10^9$ particles on multi-GPU nodes encounters a sharp memory-and-interconnect wall. Storing particle coordinates, velocities, accelerations, time-bin states, radix octree geometry, multipoles, work pools, and long-range PM buffers for $10^9+$ particles requires $\approx 140\,\mathrm{GB}$ of live state, exceeding the $48\,\mathrm{GiB}$ capacity of an individual workstation or cluster GPU such as the NVIDIA RTX A6000. Single-GPU time-slicing or host-streaming engines resolve capacity by swapping particle slabs over PCIe every timestep, which exposes host-device memory bandwidth bottlenecks whenever active particle fractions or neighbor ghost traffic grow. Conversely, naive multi-GPU distribution can introduce tree-walk divergence, dynamic memory exhaustion under late-time gravitational clustering ($z \rightarrow 0$), or peer-to-peer (P2P) hangs across non-uniform PCIe root-complex topologies.

In this work, we eliminate host particle swapping entirely at $1024^3$ while scaling across multiple GPUs inside a single node. Our solver, \texttt{gravityCU} (\texttt{multi\_gpu\_slab.cu} and \texttt{grav\_domain\_3d\_engine.cuh}), supports both cyclic 12-slab partitioning and dynamic $12\times 12\times 12$ 3D Hilbert/ORB domain decomposition across $G \in \{1, 2, 3, 4, 6\}$ GPUs, keeping every particle, persistent lazy octree, and distributed PM buffer resident inside device VRAM throughout cosmological integration, complemented by an inline/standalone Full-GPU FoF, Catalog, and Spherical Overdensity module (\texttt{grav\_cuda\_fof.cu}). Crucially, we preserve exact algorithmic parity with \texttt{SWIFT}'s cosmology, Plummer force softening, dynamic power-of-two PM timestep hierarchy, and leapfrog kick-drift-kick integrator, enabling one-to-one validation against archival 256-core CPU \texttt{SWIFT}, \texttt{GADGET-4}, \texttt{PKDGRAV3}, and reference-independent analytical physics tests.

The remainder of this paper is organized as follows. Section~\ref{sec:arch} describes the multi-GPU VRAM-resident architecture, cyclic and 3D Hilbert load balancing, NVLink vs.\ PCIe \texttt{NODE} ghost exchange, and resident cuFFT PM integration. Section~\ref{sec:bench} specifies our benchmark cosmology, archival CPU reference dataset, and accuracy metrics. Section~\ref{sec:perf} reports Gate~1 ($256^3$) strong scaling, kernel/3D optimizations and Full-GPU FoF (Figs.~\ref{fig:opt}--\ref{fig:acc}), interconnect ablations, and $1024^3$ runtime and VRAM balance. Section~\ref{sec:valid} presents end-to-end $1024^3$ accuracy validation alongside reference-independent analytical and conservation gates (Zel'dovich collapse, exact Ewald $\theta$-sweep, Plummer, and Layzer--Irvine). Section~\ref{sec:history} calibrates a six-epoch $1024^3$ full-history scaling model and presents a controlled four-code wall-time, NVML GPU energy, and normalized-throughput comparison (\texttt{SWIFT}, \texttt{GADGET-4}, \texttt{PKDGRAV3}, and \texttt{gravityCU}, Fig.~\ref{fig:crosscode}), before Section~\ref{sec:conc} concludes with software availability.

\section{Multi-GPU VRAM-Resident 12-Slab Architecture}\label{sec:arch}

\figwide{fig1_architecture.pdf}{Architecture of the multi-GPU VRAM-resident 12-slab TreePM engine on the $4 \times \text{RTX A6000}$ node \texttt{a9}. Twelve periodic $Y$-slabs (3 slabs per GPU) remain continuously resident in GPU VRAM with zero host particle swaps across timesteps. Intra-island ghost trees and slab migrants transfer directly over NVLink4 (\texttt{GPUs 0$\leftrightarrow$1} and \texttt{2$\leftrightarrow$3}, \texttt{can\_peer=1}), whereas cross-bridge exchanges between \texttt{GPU 1} and \texttt{GPU 2} (and periodic wrap \texttt{0$\leftrightarrow$11}) traverse pinned host staging (\texttt{GRAV\_P2P\_MODE=nvlink}), bypassing the $1024^3$ \texttt{NODE} PCIe P2P deadlock observed under \texttt{GRAV\_P2P\_MODE=all}.}{fig:arch}

\subsection{Cyclic Slab Domain Decomposition and Memory Plan}

To fit $1024^3$ particles cleanly across $4 \times 48\,\mathrm{GiB}$ GPUs without exceeding any device's $50.90\,\mathrm{GB}$ ($49{,}140\,\mathrm{MiB}$) ceiling even under strong nonlinear clustering at $z=0$, \texttt{multi\_gpu\_slab.cu} divides the periodic box along one Cartesian axis into $N_s = 12$ uniform slabs of width $L/12$, assigning $12/G$ contiguous slabs per GPU. Before allocating device buffers, Pass~1 streams the binary initial-state coordinates in $16\mathrm{M}$-particle chunks, builds slab histograms along all three coordinate axes ($x, y, z$), and searches all $3 \times 12$ cyclic \texttt{(axis, offset)} pairs to minimize the maximum particle count on any single GPU:
\begin{equation}
  (a^*, o^*) = \arg\min_{(a,o)} \max_{0 \le g < G}
  \sum_{\ell=0}^{N_s/G - 1} N_{a, (o + g N_s/G + \ell)\bmod 12},
\end{equation}
breaking ties by minimum peak single-slab count without modifying particle coordinates. On our $L = 25\,h^{-1}\mathrm{Mpc}$ cosmological box at both $z = 2.033$ and $z = 0.063$, Pass~1 automatically selects \texttt{slab\_axis = 1} ($Y$-axis) and \texttt{slab\_offset = 2}.

Simultaneously, Pass~1 detects uniform particle mass (\texttt{uniform\_mass = true} for dark-matter-only runs), eliding per-particle mass arrays from VRAM. Each slab allocates persistent device arrays for particle positions (\texttt{double3}), velocities (\texttt{float3}), old accelerations, time-bin tags, mesh accelerations (\texttt{float3}), sorted particle permutations, radix-tree geometry, and multipoles, sized with a conservative growth factor (\texttt{1.04} slab capacity slack for $N > 5 \times 10^7$ plus dedicated node and CUB sort-scratch pools). Thread-safe phase-pool guards (\texttt{omp critical(phase\_pool\_alloc)}) prevent concurrent OpenMP host threads from racing during per-GPU workspace initialization. Under this plan at $1024^3$ across $G=4$ GPUs, total allocated device memory is $139.69\,\mathrm{GB}$ at $z = 1.758$ (minimum free VRAM $12.31\,\mathrm{GB}$) and $142.40\,\mathrm{GB}$ at $z = 0$ (minimum free VRAM $4.86\,\mathrm{GB}$ on the densest cluster-bearing GPU~2 hosting $348.2\mathrm{M}$ particles), achieving zero host swaps over hundreds of timesteps.

\subsection{Resident Octree, Ghost Exchange, and Interconnect Topology}

Within each timestep, OpenMP host threads orchestrate the 4 GPUs through six pipelined phases (Fig.~\ref{fig:arch}):
\begin{enumerate}
  \item \textit{Drift and Dual-Stream Slab Migration}: Active and drifting particles advance on device; particles crossing a slab boundary pack into typed \texttt{DevMigrant} packets and stream asynchronously along paired \texttt{d2h\_left\_stream} and \texttt{d2h\_right\_stream} paths to adjacent slabs (wrapping periodically between slab~11 and slab~0).
  \item \textit{Resident $128^3$ cuFFT PM Solve}: Whenever the integer timeline crosses a PM boundary, each GPU deposits its local slabs onto a resident $128^3$ \texttt{double} mesh via GPU CIC (\texttt{grav\_cuda\_pm.cu} / \texttt{grav\_cuda\_pm\_dist.cu}), reduces partial meshes via host or peer-staged distributed \texttt{cuFFT} 3D Poisson solve with \texttt{SWIFT}'s exact truncated Green's function, and scatters interpolated long-range forces back to device particles.
  \item \textit{Per-Slab Radix Octree and Multipole Construction}: Each slab sorts active keys in device scratch memory, builds a $14$-level radix octree (\texttt{tree\_leaf\_limit = 64}), refits displaced cell centers in their root periodic image, and accumulates $P_2M$ multipoles up to $144$ top-level slab roots.
  \item \textit{Neighborhood Ghost-Tree Exchange}: Because the Barnes--Hut opening angle $\theta = 0.7$ and TreePM short-range cutoff $r_s$ allow direct or multipole interactions across adjacent slab boundaries, each slab imports the compact \texttt{geometry}, \texttt{multipoles}, and \texttt{force\_nodes} buffers of its $\pm 2$ neighboring slabs.
  \item \textit{Warp-Level Tree-FMM Force Integration}: Custom \texttt{tree\_force\_warp\_slab} and zero-spill split-force kernels execute warp-cooperative M2L, P2P, and fast exponential short-range split kernels (\texttt{GRAV\_FAST\_M2P\_EXP}, \texttt{warp\_m2l=1}, \texttt{fast\_split=1}) with zero host involvement.
  \item \textit{Leapfrog Kick Synchronization}: \texttt{integrate\_slab} closes completed \texttt{kick2} and PM kicks using \texttt{finish[old]}, updates per-particle adaptive power-of-two time bins (\texttt{timeline} $= 2^{57}$), and opens new \texttt{kick1} steps except at the final step \texttt{ti == timeline}.
\end{enumerate}

A critical systems finding emerged when deploying this ghost-exchange pattern on dual-socket AMD EPYC GPU nodes such as node~\texttt{a9}. Hardware topology inspection via \texttt{nvidia-smi topo -m} reveals that GPUs~\texttt{0..3} share NUMA domain~0 but are split across two PCIe root complexes: GPUs~\texttt{0$\leftrightarrow$1} and GPUs~\texttt{2$\leftrightarrow$3} are connected by bonded 4-lane NVLink4 bridges (\texttt{NV4}), whereas pairs crossing the bridge (\texttt{0,1}$\leftrightarrow$\texttt{2,3}) communicate over PCIe \texttt{NODE} interconnects without hardware NVLink atomics. When \texttt{GRAV\_P2P\_MODE=all} enables direct CUDA \texttt{cudaDeviceEnablePeerAccess} across \texttt{NODE} links (\texttt{can\_peer[i][j] = 1} for all pairs), $1024^3$ execution deadlocks on Step~0 with GPUs~0 and~2 spinning at $100\%$ SM utilization ($140\,\mathrm{W}$) and GPUs~1 and~3 stalled idle ($0\%$, $21\text{--}24\,\mathrm{W}$). Setting \texttt{GRAV\_P2P\_MODE=nvlink} (default) gates direct \texttt{cudaDeviceEnablePeerAccess} strictly on \texttt{cudaDevP2PAttrNativeAtomicSupported == 1}, routing intra-pair transfers (\texttt{0$\leftrightarrow$1}, \texttt{2$\leftrightarrow$3}) via zero-copy NVLink and inter-island ghost/migrant transfers (\texttt{1$\leftrightarrow$2}, periodic \texttt{0$\leftrightarrow$3}) cleanly through dual pinned-host staging buffers.

\subsection{Fast Single-Pass C CIC Deposition for Billion-Particle Analysis}

Analyzing billion-particle HDF5 snapshots in pure Python via repeated \texttt{numpy.add.at} scatter operations across $1{,}024$ one-million-particle chunks ($\sim 5.7 \times 10^4$ Python kernel dispatches onto $256^3$ double grids) consumed $42$--$52\,\mathrm{min}$ on host CPUs and triggered a $2.2\,\mathrm{h}$ idle-GPU Slurm termination in an initial mid-redshift run (\texttt{Slurm 489937}). We resolved this analysis bottleneck by implementing a single-pass deterministic C99 CIC kernel (\texttt{cic\_deposit.c}) invoked via \texttt{ctypes} from \texttt{eval\_1024\_snapshot.py}. Verified against \texttt{numpy.add.at} on $1{,}048{,}576$ particles including periodic wrap boundaries, the C kernel achieves $\max|\Delta \rho| = 4.44 \times 10^{-16}$ and $\max|\Delta \mathbf{p}| = 2.27 \times 10^{-13}$ (particle mass conservation error $0.0$) while accelerating CIC deposition by $8.3\times$, cutting full evaluation wall time on snapshot~\texttt{0045} to $27.9\,\mathrm{min}$ (dominated by serial \texttt{gzip-1} HDF5 decompression of $\sim 72\,\mathrm{GB}$ input).

\section{Benchmark Suite and Validation Methodology}\label{sec:bench}

All benchmarks solve pure collisionless dark-matter gravity in a periodic comoving box of $L = 25\,h^{-1}\,\mathrm{Mpc}$ ($37.1140\,\mathrm{Mpc}$) under a Planck~2018~\citep{Planck2018VI} flat $\Lambda$CDM cosmology ($h = 0.6736$, $\Omega_m = 0.3153$, $\Omega_\Lambda = 0.6847$, $\Omega_{\rm cdm} = 0.266$), matching an archival 256-core CHPC CPU \texttt{SWIFT} production simulation (\texttt{PL\_25\_1024}) initialized at $z = 200$ via \texttt{monofonIC}~\citep{Michaux2021monofonIC} and evolved through $20{,}194$ clean steps (\texttt{0..20193}) to $z = 0$ with $56$ snapshots (\texttt{0000}--\texttt{0056}). Key runtime parameters match upstream \texttt{SWIFT}: Plummer softening $\epsilon = 2.1746\,\mathrm{kpc}$ ($1.465\,h^{-1}\,\mathrm{kpc}$ comoving), Barnes--Hut opening angle $\theta = 0.7$, dynamic power-of-two PM spans with \texttt{16} short substeps per PM interval, and gravitational constant $G = 43.0092$.

Our validation ladder consists of three stages executed on node~\texttt{a9} ($4 \times \text{NVIDIA RTX A6000}$, dual-socket AMD EPYC host, CUDA~12.6):
\begin{enumerate}
  \item \textit{Gate~1 ($256^3 = 16{,}777{,}216$ particles, $a = 0.855 \rightarrow 0.920$)}: Verifies algorithmic invariance and strong scaling across $G = 1, 2, 4$ GPUs (\texttt{Slurm 489828}, $320$ steps, $17$ PM updates) under identical parameters.
  \item \textit{Late-Redshift $1024^3$ End-to-End ($0055 \rightarrow 0056$, $z = 0.063 \rightarrow 0$, $a = 0.94059 \rightarrow 1.00000$)}: Evolves $1.074 \times 10^9$ particles from archival snapshot \texttt{PL\_25\_1024\_0055.hdf5} through $513$ GPU steps ($33$ PM updates) across $1\times$ RTX A6000, $4\times$ single-NVLink (\texttt{Slurm 489830}), $4\times$ \texttt{GRAV\_P2P\_MODE=all} probe (\texttt{Slurm 489878}), and $4\times$ dual-NVLink (\texttt{Slurm 489882}).
  \item \textit{Mid-Redshift $1024^3$ End-to-End ($0044 \rightarrow 0045$, $z = 2.033 \rightarrow 1.758$, $a = 0.32967 \rightarrow 0.36264$)}: Evolves $1.074 \times 10^9$ particles from \texttt{PL\_25\_1024\_0044.hdf5} through $910$ GPU steps ($50$ PM updates) on $4\times$ dual-NVLink GPUs both solo (\texttt{Slurm 489937}) and co-scheduled on a shared node (\texttt{Slurm 489982}), complemented by a $17$-step $1\times$ RTX A6000 timing probe ($a = 0.32967 \rightarrow 0.33050$).
\end{enumerate}

For every $1024^3$ end-to-end validation, GPU binaries are converted back to HDF5, clustered with standalone native 64/96-thread CPU \texttt{SWIFT fof}~\citep{Davis1985FOF} using identical configuration parameters, and evaluated against archival CHPC snapshots (\texttt{PL\_25\_1024\_0056/0045.hdf5} and \texttt{fof\_output\_0056/0045.hdf5}) across four metrics: (i) $16{,}388$ ID-matched particle position and velocity errors (plus \texttt{GRAV\_FORCE3\_OUT} short/long force diagnostics for kick-convention alignment), (ii) CIC-deconvolved $256^3$ matter density power spectra $P_\delta(k)$ over $k \in [0.3, 16]\,h\,\mathrm{Mpc}^{-1}$ ($63$ spherical $k$-bins), (iii) velocity-momentum divergence power spectra $P_\theta(k)$ before and after short-kick synchronization, and (iv) cumulative FoF halo mass functions over $N_{\min} \in [32, 262{,}144]$. Device power and memory telemetry are sampled at $200\,\mathrm{ms}$--$5\,\mathrm{s}$ via \texttt{nvidia-smi} filtered strictly to the allocated GPU indices.

\section{Performance, Strong Scaling, and Interconnect Topology}\label{sec:perf}

\subsection{Gate 1 Strong Scaling and Multi-GPU Invariance (\texorpdfstring{$256^3$}{256 cubed})}

Table~\ref{tab:gate1} summarizes Gate~1 on $256^3$ particles across $G = 1, 2, 4$ GPUs (\texttt{CUDA\_VISIBLE\_DEVICES=1,2,3,4}, where only GPUs~\texttt{2$\leftrightarrow$3} share NVLink). Scaling from $1$ to $2$ GPUs uses pinned-host staging on both boundaries (\texttt{can\_peer = [[1,0],[0,1]]}) and trims runtime from $72.80\,\mathrm{s}$ to $61.48\,\mathrm{s}$ ($1.18\times$). Enabling $G = 4$ adds NVLink peer access on the central boundary (\texttt{2$\leftrightarrow$3}) and drops runtime to $30.78\,\mathrm{s}$ ($2.36\times$ vs.\ $G=1$). Crucially, all three GPU counts preserve bit-level scientific invariance relative to the CPU reference (\texttt{pos\_rms = 0.022825 kpc/h}, $P_\delta$ max deviation $1.057 \times 10^{-5}$, and kick-aligned $P_\theta$ max deviation $1.946 \times 10^{-4}$ across all $G$).

\begin{table}[!tbp]
\caption{Gate 1 ($256^3$, $a = 0.855 \rightarrow 0.920$, 320 steps, 17 PM) strong scaling and invariance on node \texttt{a9} (\texttt{Slurm 489828}).}\label{tab:gate1}
\scriptsize
\setlength{\tabcolsep}{2.5pt}
\begin{ruledtabular}
\begin{tabular}{lccccc}
GPUs & Topology & Evo (s) & vs.\ $G=1$ & Pos RMS & $P_\delta$ max dev \\
\hline
$G=1$ & Single GPU & $72.80$ & $1.00\times$ & $0.02283$ & $1.057\times 10^{-5}$ \\
$G=2$ & PCIe bounce & $61.48$ & $1.18\times$ & $0.02283$ & $1.057\times 10^{-5}$ \\
$G=4$ & Single-NVL & $30.78$ & $2.36\times$ & $0.02283$ & $1.057\times 10^{-5}$ \\
\end{tabular}
\end{ruledtabular}
\end{table}

\subsection{Billion-Particle Timing, Interconnect Effects, and VRAM Headroom}

\figwide{fig2_scaling_and_vram.pdf}{(a) Pure gravity evolution wall-clock time and speedup vs.\ the 256-core archival CHPC CPU run for $1024^3$ segments at $z = 0.063 \rightarrow 0$ ($513$ GPU steps) and $z = 2.033 \rightarrow 1.758$ ($910$ GPU steps); the hatched bar denotes the $1\times$ A6000 $0044 \rightarrow 0045$ projection calibrated from a $17$-step mid-$z$ probe ($265.68\,\mathrm{s}$). (b) Per-GPU device memory allocation and particle counts on dual-NVLink $4 \times \text{RTX A6000}$ at $z = 1.758$ ($12.31\,\mathrm{GB}$ minimum free headroom) and $z = 0$ ($4.86\,\mathrm{GB}$ minimum free headroom).}{fig:scaling}

Table~\ref{tab:perf1024} and Fig.~\ref{fig:scaling} quantify billion-particle ($1024^3$) scaling across interconnects and redshifts. At $z = 0.063 \rightarrow 0$, moving from $1\times$ RTX A6000 (\texttt{slab\_resident.cu}, $10{,}540.44\,\mathrm{s}$) to $4\times$ RTX A6000 with a single active NVLink pair (\texttt{1,2,3,4}, \texttt{Slurm 489830}) cuts runtime to $5{,}719.90\,\mathrm{s}$ ($1.84\times$ multi-GPU scaling; $5.57\times$ vs.\ 256-core CPU). Re-targeting the job to dual-NVLink pairs (\texttt{0,1,2,3}, \texttt{Slurm 489882}) eliminates the extra host-bounce boundary between GPUs~0 and~1, cutting runtime by a further $1.30\times$ down to \textbf{$4{,}388.62\,\mathrm{s}$ ($73\,\mathrm{m}\,09\,\mathrm{s}$)}---a \textbf{$2.40\times$ speedup over $1\times$ A6000} and \textbf{$7.26\times$ speedup over 256-core CPU \texttt{SWIFT}} ($31{,}878.77\,\mathrm{s}$ across $757$ CPU steps). Conversely, attempting direct CUDA P2P over PCIe \texttt{NODE} bridges (\texttt{SWIFT\_P2P\_MODE=all}, \texttt{Slurm 489878}) hangs indefinitely at Step~0, validating native-atomic NVLink gating as essential on dual-root EPYC hardware.

At mid-redshift ($z = 2.033 \rightarrow 1.758$), gravitational clustering is weaker and particle distribution across the 12 $Y$-slabs is significantly more uniform: GPU~particle loads range over $[205.0, 279.0, 293.1, 296.6]\,\mathrm{M}$ versus $[137.4, 269.5, 348.2, 318.7]\,\mathrm{M}$ at $z=0$. Consequently, minimum per-GPU free VRAM rises from $4.86\,\mathrm{GB}$ ($z=0$) to \textbf{$12.31\,\mathrm{GB}$} ($z=1.758$), and pure evolution on an idle \texttt{a9} node (\texttt{Slurm 489937}) finishes $910$ steps ($50$ PM updates) in \textbf{$3{,}510.98\,\mathrm{s}$ ($58\,\mathrm{m}\,31\,\mathrm{s}$)}. Against the 256-core archival CPU run ($589$ steps totaling $33{,}863.74\,\mathrm{s}$, or $32{,}538.37\,\mathrm{s}$ pure gravity excluding Step~8392 inline snapshot/FoF dump overhead), dual-NVLink $4\times$ A6000 delivers \textbf{$9.27\times$ pure-evolution speedup} ($9.64\times$ against total CPU step time). Even when node~\texttt{a9} is shared concurrently with another user's GPU job occupying GPU~4 and 128 host cores (\texttt{Slurm 489982}), shared-PCIe contention slows evolution by only $24\%$ to $4{,}362.17\,\mathrm{s}$ ($7.46\times$ pure / $7.76\times$ total vs.\ CPU), and allocated GPU power draw drops correspondingly ($637.14\,\mathrm{W}$ vs.\ $735.87\,\mathrm{W}$ solo), keeping GPU energy consumption nearly invariant ($0.772\,\mathrm{kWh}$ vs.\ $0.718\,\mathrm{kWh}$).

\begin{table*}[!tbp]
\caption{Billion-particle ($1024^3$) execution performance, minimum free VRAM headroom, and allocated GPU telemetry power on $1\times$ and $4\times$ RTX A6000 vs.\ 256-core CPU \texttt{SWIFT}.}\label{tab:perf1024}
\footnotesize
\setlength{\tabcolsep}{3pt}
\begin{ruledtabular}
\begin{tabular}{llcccccc}
Epoch / Run & Config & Steps (PM) & Evo (s) & vs.\ CPU & vs.\ $1\times$ GPU & Free (GB) & Power (W) \\
\hline
\multicolumn{8}{c}{\textit{Late Redshift: $0055 \rightarrow 0056$ ($z = 0.063 \rightarrow 0$, CPU: $757$ steps / $31{,}878.77\,\mathrm{s}$)}} \\
\texttt{slab-1024-0056} & $1\times$ A6000 resident & 513 (33) & $10{,}540.44$ & $3.02\times$ & $1.00\times$ & $17.34$ & $217.41$ \\
\texttt{Slurm 489830} & $4\times$ A6000 1-NVLink & 513 (33) & $5{,}719.90$ & $5.57\times$ & $1.84\times$ & $4.86$ & --- \\
\texttt{Slurm 489878} & $4\times$ A6000 \texttt{p2p=all} & --- & \multicolumn{3}{c}{\textit{Stalled at Step~0 (PCIe NODE P2P hang)}} & $4.86$ & --- \\
\texttt{Slurm 489882} & $4\times$ A6000 Dual-NVL & 513 (33) & $\mathbf{4{,}388.62}$ & $\mathbf{7.26\times}$ & $\mathbf{2.40\times}$ & $4.86$ & $657.12$ \\
\hline
\multicolumn{8}{c}{\textit{Mid Redshift: $0044 \rightarrow 0045$ ($z = 2.033 \rightarrow 1.758$, CPU: $589$ steps / $32{,}538.37\,\mathrm{s}$ pure, $33{,}863.74\,\mathrm{s}$ tot)}} \\
\texttt{slab1-probe} & $1\times$ A6000 (17-step) & 17 (2) & $265.68^\dagger$ & $3.35\times^\dagger$ & $1.00\times$ & $24.89$ & --- \\
\texttt{Slurm 489937} & $4\times$ A6000 Dual-NVL solo & 910 (50) & $\mathbf{3{,}510.98}$ & $\mathbf{9.27\times}$ & $\approx 2.77\times^\dagger$ & $\mathbf{12.31}$ & $735.87$ \\
\texttt{Slurm 489982} & $4\times$ A6000 Dual-NVL shr. & 910 (50) & $4{,}362.17$ & $7.46\times$ & $\approx 2.23\times^\dagger$ & $12.31$ & $637.14$ \\
\end{tabular}
\end{ruledtabular}
\begin{flushleft}
\footnotesize $^\dagger$The 17-step mid-$z$ probe ($a = 0.32967 \rightarrow 0.33050$, mean PM $59.85\,\mathrm{s}$, non-PM fit $4.30 + 69.83\,f_{\rm act}\,\mathrm{s}$) projects $9{,}710.15\,\mathrm{s}$ over the 910-step workload.
\end{flushleft}
\end{table*}

\subsection{Kernel Optimizations, 3D Domain Decomposition, and Full-GPU FoF}

\figwide{fig5_opt_speedup_vs_cpu.pdf}{(a) $256^3$ ($a = 0.855 \rightarrow 0.920$, $320$ steps) evolution progression from initial $1\times$ A6000 ($126.8\,\mathrm{s}$) to $4\times$ A6000 3D Hilbert Full ($17.53\,\mathrm{s}$, $45.44\times$ vs.\ 256-thread CPU \texttt{SWIFT}, or $50.6\times$ excluding export I/O). (b) Phase breakdown showing tree-build overhead reduced by $96.2\%$ ($10.09\,\mathrm{s} \rightarrow 0.38\,\mathrm{s}$). (c) Multi-scale gravity and FoF/SO speedups. (d) $1024^3$ measured windows and full-history projections.}{fig:opt}

\figwide{fig6_accuracy_vs_cpu.pdf}{Scientific accuracy after all kernel, 3D Hilbert, and Full-GPU FoF optimizations. (a)--(b) $256^3$ matter and kick-aligned momentum power spectra remain within $0.00106\%$ (eager) / $0.0387\%$ (lazy) and $0.0198\%$--$0.0578\%$, respectively. (c) Error ladder spanning double-precision Catalog/SO parity ($\le 1.16\times 10^{-13}$ and $\le 3.67\times 10^{-15}$), $256^3$/$1024^3$ spectra, and $0/16.78\mathrm{M}$ FoF halo-ID mismatches. (d) $1024^3$ power spectra at $z=0$ and $z=1.758$.}{fig:acc}

Beyond the 12-slab baseline, Priorities~1--5 refactor the solver in \texttt{grav\_domain\_3d\_engine.cuh}, \texttt{grav\_cuda\_pm\_dist.cu}, and \texttt{grav\_cuda\_fof.cu} into five composable optimizations evaluated at $256^3$ ($N = 16{,}777{,}216$, Figs.~\ref{fig:opt}--\ref{fig:acc}):
(i) \textit{Lazy Octree Refit} (\texttt{GRAV\_LAZY\_TREE=1}) reuses persistent tree topology across active substeps, cutting $1\times$-GPU tree build time by $84\%$ ($9.96\,\mathrm{s} \rightarrow 1.58\,\mathrm{s}$);
(ii) \textit{Zero-Spill Split Force} (\texttt{GRAV\_SPLIT\_FORCE=1}) decomposes tree traversal, M2P accumulation, grouped M2L, and warp-tiled P2P into register-lean kernels with zero local-memory spill;
(iii) \textit{Pure-CUDA Distributed PM} (\texttt{GRAV\_PM\_DIST=1}) parallelizes $128^3$ cuFFT Poisson solves across GPUs using NVLink P2P + pinned bounce buffers;
(iv) \textit{Dynamic 3D $12\times 12\times 12$ Hilbert Domain Decomposition} (\texttt{GRAV\_DOMAIN\_MODE=3d\_hilbert}) achieves interior-first overlap and lowers $4\times$ A6000 evolution time to \textbf{$17.53\,\mathrm{s}$} ($15.73\,\mathrm{s}$ compute+PM, \textbf{$45.44\times$} vs.\ $796.36\,\mathrm{s}$ 256-thread CPU \texttt{SWIFT}, and $2.99\times$ scaling over $1\times$ GPU $52.44\,\mathrm{s}$), reducing tree overhead to $0.38\,\mathrm{s}$ ($-96.2\%$ vs.\ $10.09\,\mathrm{s}$ pre-opt) while keeping $P_\delta(k)$ within $0.00106\%$ (eager) and $0.0387\%$ (lazy); and
(v) \textit{Full-GPU FoF, Catalog, and SO} (\texttt{grav\_cuda\_fof.cu}) runs zero-upload 3D/slabbed Union--Find FoF directly on resident device coordinates with $0 / 16.78\mathrm{M}$ halo-ID mismatches across $S \in \{4, 8\}$ and $G \in \{1, 2, 4\}$, accelerating halo catalog reduction by \textbf{$12.48\times$} ($\max\text{rel} \le 1.16\times 10^{-13}$) and segmented-sort spherical overdensity ($M_{200c}/R_{200c}$) by \textbf{$15.33\times$} uniform / \textbf{$5.37\times$} non-uniform ($\max\text{rel} \le 3.67\times 10^{-15}$) relative to 8-thread OpenMP CPU baselines.

\section{Multi-Redshift Scientific Validation}\label{sec:valid}

\figwide{fig3_scientific_validation.pdf}{Billion-particle ($1024^3$) scientific accuracy of the $4 \times \text{RTX A6000}$ Dual-NVLink solver against archival 256-core CPU \texttt{SWIFT} at $z=0$ (\texttt{0056}) and $z=1.758$ (\texttt{0045}). (a) Matter power spectrum fractional deviation $(P_\delta^{\rm GPU}/P_\delta^{\rm CPU} - 1)$ across $k \in [0.3, 16]\,h\,\mathrm{Mpc}^{-1}$ remains within $0.0045\%$ ($z=0$) and $0.0102\%$ ($z=1.758$). (b) Momentum divergence spectrum $P_\theta(k)$: raw $z=1.758$ velocities match directly ($0.0111\%$), whereas end-of-run $z=0$ velocities exhibit a $-2\mathbf{s}$ short-force kick convention offset ($0.702\%$ raw) whose removal collapses the residual $17.3\times$ to $0.0405\%$. (c) Cumulative native \texttt{SWIFT} FoF halo counts agree within $\pm 0.32\%$ ($z=0$) and $\pm 0.18\%$ ($z=1.758$, $N_{\min} \le 65{,}536$), well inside Poisson $\pm N_{\rm halo}^{-1/2}$ bounds.}{fig:valid}

\subsection{Particle Accuracy and Matter Power Spectra}

As summarized in Table~\ref{tab:acc} and Fig.~\ref{fig:valid}(a), every $1024^3$ validation finishes with \texttt{bad\_count = 0}. Across $16{,}388$ ID-matched particles (\texttt{ID mod 65521 == 0}), comoving position errors remain deep in the sub-softening regime: at $z=0$ (\texttt{0056}), $\mathrm{RMS} = 0.05051\,h^{-1}\mathrm{kpc}$ (\texttt{p50 = 0.01019 kpc/h}, $< 1\%$ of the $1.465\,h^{-1}\mathrm{kpc}$ softening length); at $z = 1.758$ (\texttt{0045}), $\mathrm{RMS} = 0.09717\,h^{-1}\mathrm{kpc}$ with an even tighter median of \texttt{0.00775 kpc/h}. Comparing single-GPU (\texttt{slab-1024-0056}) against $4\times$-GPU (\texttt{Slurm 489882}) at $z=0$ confirms that multi-GPU slab decomposition introduces less than $10^{-5}$ relative difference in RMS error ($0.050507$ vs.\ $0.050513\,h^{-1}\mathrm{kpc}$).

On the deconvolved $256^3$ FFT density field over $k \in [0.3, 16]\,h\,\mathrm{Mpc}^{-1}$ ($63$ bins), the matter power spectrum $P_\delta(k)$ matches archival CPU \texttt{SWIFT} to $\max|P_\delta^{\rm GPU}/P_\delta^{\rm CPU} - 1| = \mathbf{0.00453\%}$ ($\mathrm{RMS} = 0.00288\%$, $\min r = 0.999999973$) at $z=0$ and $\mathbf{0.01019\%}$ ($\mathrm{RMS} = 0.00610\%$, $\min r = 0.999999950$) at $z = 1.758$. On cosmological linear-to-quasilinear scales ($0.3\text{--}1\,h\,\mathrm{Mpc}^{-1}$), deviations do not exceed $0.00017\%$ at either redshift, outperforming standard precision-cosmology accuracy requirements ($\sim 0.1\text{--}1\%$) by more than an order of magnitude.

\begin{table}[!tbp]
\caption{End-to-end $1024^3$ accuracy on $4\times$ RTX A6000 Dual-NVLink relative to archival 256-core CPU \texttt{SWIFT} ($0.3\text{--}16\,h\,\mathrm{Mpc}^{-1}$).}\label{tab:acc}
\scriptsize
\setlength{\tabcolsep}{1.1pt}
\begin{ruledtabular}
\begin{tabular}{lcc}
Metric & \texttt{0056} ($z=0$) & \texttt{0045} ($z=1.758$) \\
\hline
Pos RMS / p50 (kpc/$h$) & $0.0505$ / $0.0102$ & $0.0972$ / $0.0078$ \\
$P_\delta(k)$ max dev / RMS & $\mathbf{0.00453\%}$ / $0.0029\%$ & $\mathbf{0.01019\%}$ / $0.0061\%$ \\
$P_\delta(k)$ min $r$ & $0.999999973$ & $0.999999950$ \\
Raw vel RMS / p50 & $2.666$ / $0.996$ & $\mathbf{0.623}$ / $\mathbf{0.0186}$ \\
Sync vel RMS / p50 & $\mathbf{1.187}$ / $\mathbf{0.336}$ & \textit{(already sync.)} \\
Raw $P_\theta(k)$ max dev & $0.7018\%$ ($0.99976$) & $\mathbf{0.01112\%}$ ($0.9999999$) \\
Aligned $P_\theta(k)$ dev & $\mathbf{0.0405\%}$ ($0.9999996$) & \textit{(over-shift $0.569\%$)} \\
FoF $\#1$ halo & $34.804\mathrm{M}$ vs.\ $34.805\mathrm{M}$ & $6.566\mathrm{M}$ vs.\ $6.563\mathrm{M}$ \\
FoF $32\text{--}65\text{k}$ diff & $-0.32\%\,..\,{+}0.03\%$ & $-0.18\%\,..\,{+}0.02\%$ \\
FoF $N\!\ge\!262{,}144$ & $198$ vs.\ $198$ ($0.0\%$) & $146$ vs.\ $147$ ($-0.68\%$) \\
\end{tabular}
\end{ruledtabular}
\end{table}

\subsection{Velocity Kick Synchronization Closure Across Snapshots}

Examining the momentum divergence spectrum $P_\theta(k)$ ($\theta = \nabla \cdot [(1+\delta)\mathbf{v}]$) across both redshifts reveals an instructive contrast rooted in leapfrog snapshot timing (Fig.~\ref{fig:valid}(b)). In \texttt{integrate\_slab}, when the timeline reaches the run endpoint \texttt{ti == timeline}, the GPU closes both short-range (\texttt{kick2}) and PM kicks (\texttt{mesh\_kick}) while suppressing opening \texttt{kick1} (\texttt{if (ti != timeline)}), exporting velocities strictly synchronized at $a_{\rm end}$.

At the intermediate archival snapshot \texttt{0045} (CHPC CPU step \texttt{8392}, property bitmask \texttt{Props = 145 = Rebuild + Snapshot + FOF}), archival CPU velocities match the GPU synchronized export directly: raw velocity differences show a tiny median of \textbf{$0.0186\,\mathrm{km/s}$} ($\mathrm{RMS} = 0.6226\,\mathrm{km/s}$) and raw $P_\theta(k)$ agrees to \textbf{$0.01112\%$} ($\min r = 0.999999879$); artificially adding $+2\mathbf{s}$ in post-processing un-synchronizes the velocities ($P_\theta$ degrades to $0.5694\%$). By contrast, at the final end-of-simulation snapshot \texttt{0056} (CHPC step \texttt{20193}, \texttt{Props = 257 = Rebuild + mesh}, $g$-updates $= 1{,}073{,}741{,}824$), raw CPU--GPU velocity residuals regress linearly onto the per-particle final short-range kick $\mathbf{s}$ and long-range kick $\mathbf{l}$ with two-component slope $[-1.9197, -0.0219]$, explaining \textbf{$80.30\%$} of raw $\Delta v$ variance ($2.6657 \rightarrow 1.1870\,\mathrm{km/s}$) and collapsing $P_\theta(k)$ error \textbf{$17.3\times$} from $0.7018\%$ down to \textbf{$0.0405\%$} ($\min r = 0.999999614$). Taken together, both snapshots confirm $\sim 0.01\text{--}0.04\%$ intrinsic velocity-divergence agreement once endpoint kick state is matched.

\subsection{Native SWIFT FoF Halo Mass Function Alignment}

Running upstream native \texttt{SWIFT fof} directly on the $z=0$ ($64$ threads) and $z=1.758$ ($96$ threads) snapshots demonstrates remarkable nonlinear halo preservation (Fig.~\ref{fig:valid}(c)). At $z=0$, the largest cluster-scale halo contains $34{,}804{,}209$ particles on $4\times$ GPU vs.\ $34{,}804{,}779$ on CPU ($-0.0016\%$ difference); cumulative halo counts across $N_{\min} \ge 32 \dots 65{,}536$ ($674{,}485$ down to $617$ halos) agree within $-0.32\%$ to $+0.026\%$, and massive halos with $N_{\min} \ge 262{,}144$ match identically ($198$ vs.\ $198$, $0.0\%$). At $z=1.758$, where halos are smaller ($6{,}566{,}460$ vs.\ $6{,}562{,}673$ particles in halo~\#1, $+0.058\%$), cumulative counts over $N_{\min} \ge 32 \dots 65{,}536$ ($871{,}028$ to $593$ halos) remain within $-0.18\%$ to $+0.019\%$, and $N_{\min} \ge 262{,}144$ differs by a single object ($146$ vs.\ $147$, $-0.68\%$), well inside Poisson $\pm N_{\rm halo}^{-1/2}$ statistical bounds ($8.2\%$).

\subsection{Reference-Independent Analytical and Conservation Gates}

To validate gravitational accuracy independently of any single external simulation code, Table~\ref{tab:physgates} compiles five self-contained analytical, exact-force, and conservation checks:
\begin{enumerate}
  \item \textit{1D Zel'dovich Plane-Wave Collapse}~\citep{Zeldovich1970}: Evolving $64^3 = 262{,}144$ particles in an Einstein--de Sitter cosmology ($\Omega_m = 1$, $L = 37.114\,\mathrm{Mpc}$, $H_0 = 67.36\,\mathrm{km\,s^{-1}\,Mpc^{-1}}$) from $a = 0.02 \rightarrow 0.50$ ($\delta_{\max} = 0.50$, $513$ steps) reproduces the exact analytical displacement ($2.953\,\mathrm{Mpc}$ peak) to $\Delta x_{\rm RMS}/L = 4.41\times 10^{-4}$ ($\Delta x_{\rm RMS} = 0.0164\,\mathrm{Mpc}$, $0.554\%$ of amplitude) and peak peculiar velocity ($281.35\,\mathrm{km/s}$) to $\Delta v_{\rm RMS}/v_{\rm peak} = 0.592\%$ ($1.665\,\mathrm{km/s}$, transverse spurious RMS $0.070\,\mathrm{km/s}$).
  \item \textit{Periodic Ewald Force Convergence vs.\ $\theta$}: On $4{,}096$ cosmological target particles ($64^3$ subsampled at $a = 0.855$), comparing \texttt{gravityCU} against direct periodic Ewald summation shows monotonic short-range median force error reduction across $\theta \in \{0.9, 0.7, 0.5, 0.3\}$ from $8.48\times 10^{-3}$ down to $1.14\times 10^{-4}$, while total force median error converges to the $128^3$ CIC PM grid floor ($1.98\times 10^{-3}$ at $\theta = 0.7$, $1.03\times 10^{-3}$ at $\theta = 0.3$).
  \item \textit{32K Isolated Equilibrium Plummer Sphere}: Integrating $32{,}768$ particles over $t = 0 \rightarrow 4.0$ preserves total energy to $|\Delta E / E_0| = 1.15\times 10^{-5}$, linear momentum drift to $6.38\times 10^{-6}$, and virial equilibrium $2T/|U| \in [0.9956, 0.9988]$.
  \item \textit{Full-History Cosmic Energy Closure}~\citep{Layzer1963}: A complete $256^3$ 3D Hilbert run from $z = 200 \rightarrow 0$ ($9{,}257$ steps, $16$ dynamic rebalances, $526.65\,\mathrm{s}$ on $4\times$ A6000) achieves Layzer--Irvine energy conservation residual $0.0307\%$, zero mass loss, momentum drift $3.72\times 10^{-5}$, $z=0$ $P_\delta(k)$ agreement $0.123\%$, and exact massive halo counts ($89/89$ for $N\ge 10^4$, $18/18$ for $N\ge 10^5$).
  \item \textit{Cross-Domain Boundary Force Parity}: Comparing $4\times$-GPU 3D Hilbert against $1\times$-GPU 1D slab forces at Step~0 confirms $\mathrm{med}|\Delta \mathbf{a}|/|\mathbf{a}| = 1.03\times 10^{-4}$ box-wide and $1.15\times 10^{-4}$ across domain boundary faces (potential median $8.88\times 10^{-6}$).
\end{enumerate}

\begin{table}[!tbp]
\caption{Reference-independent analytical, exact-Ewald, and conservation verification gates.}\label{tab:physgates}
\scriptsize
\setlength{\tabcolsep}{0.9pt}
\begin{ruledtabular}
\begin{tabular}{lll}
Gate & Metric & Result \\
\hline
Zel'dovich $64^3$ & $\Delta x_{\rm RMS}/L$; $\Delta v/v_{\rm pk}$ & $4.41\!\times\!10^{-4}$; $0.59\%$ \\
Ewald Short ($\theta{:}0.9{\to}0.3$) & Med $|\Delta\mathbf{a}_{\rm s}|/|\mathbf{a}_{\rm s}|$ & $8.48{\to}0.114\!\times\!10^{-3}$ \\
Ewald Total ($\theta{=}0.7/0.3$) & Med $|\Delta\mathbf{a}|/|\mathbf{a}|$ & $1.98$/$1.03\!\times\!10^{-3}$ \\
Plummer 32K ($t{:}0{\to}4$) & $|\Delta E/E_0|$; $2T/|U|$ & $1.15\!\times\!10^{-5}$; $0.996\text{--}0.999$ \\
$256^3$ $z{=}200{\to}0$ (4G) & L--I; $P_\delta$; $N{\ge}10^4$ & $0.031\%$; $0.12\%$; $89/89$ \\
4G 3D vs.\ 1G 1D & Face / all $|\Delta\mathbf{a}|/|\mathbf{a}|$ & $1.15$/$1.03\!\times\!10^{-4}$ \\
\end{tabular}
\end{ruledtabular}
\end{table}

\section{Full-History (\texorpdfstring{$z = 200 \rightarrow 0$}{z = 200 -> 0}) Compute and Energy Scaling}\label{sec:history}

\figwide{fig4_full_history_scaling.pdf}{(a) Six-epoch breakdown of the $20{,}194$-step archival 256-core CPU \texttt{SWIFT} run ($304.48\,\mathrm{h}$ pure gravity plus $18.18\,\mathrm{h}$ across $56$ inline Snapshot/FoF dumps) compared against clamped-interval projections for $1\times$ and $4\times$ RTX A6000 NVLink. Notice the $265$-step $z=200\rightarrow 50$ regime consumes $60.1\,\mathrm{h}$ on CPU due to uniform all-particle PM tree walks. (b) Integrated $z = 200 \rightarrow 0$ and $z = 50 \rightarrow 0$ pure-evolution wall time ($\mathrm{h}$) and whole-node energy ($\mathrm{kWh}$); error bars bracket the clamped-interval and step-physics models.}{fig:history}

To quantify end-to-end cosmological campaign savings beyond individual redshift windows, we parsed all $28{,}704$ non-comment rows of the archival CHPC timing log (\texttt{timesteps\_1024.txt}), filtering out $335$ truncated restart lines to reconstruct the complete $20{,}194$ consecutive steps (\texttt{0..20193}, $z = 200 \rightarrow 0$) summarized in Table~\ref{tab:epochs}. Separating the $56$ inline Snapshot+FoF dump steps ($18.18\,\mathrm{h}$ excess overhead above epoch step medians) yields \textbf{$304.48\,\mathrm{h}$ ($12.69\,\mathrm{days}$)} of pure CPU gravity evolution ($322.66\,\mathrm{h} = 13.44\,\mathrm{days}$ total, consuming $82{,}600$ CPU core-hours). From $z = 50 \rightarrow 0$ (steps \texttt{265..20193}, $19{,}929$ steps), CPU gravity totals $248.65\,\mathrm{h}$ ($262.54\,\mathrm{h}$ total). Notably, the first $265$ steps ($z = 200 \rightarrow 50$, just $1.3\%$ of total steps) consume $55.82\,\mathrm{h}$ of pure CPU gravity because every step triggers a full-box global PM+tree pass across all $1.074 \times 10^9$ particles over nearly unclustered cells ($738.6\,\mathrm{s}$ median step time on 256 cores), whereas uniform-grid Step-0 GPU PM steps execute in only $11.5$--$20.7\,\mathrm{s}$ on $4\times$ A6000.

\begin{table}[!tbp]
\caption{Six-epoch breakdown of the $20{,}194$-step archival 256-core CPU \texttt{PL\_25\_1024} simulation.}\label{tab:epochs}
\scriptsize
\setlength{\tabcolsep}{2.5pt}
\begin{ruledtabular}
\begin{tabular}{lcccccc}
Epoch & Steps & PM & Dumps & Pure (h) & Tot (h) & PM med (s) \\
\hline
$z = 200 \rightarrow 50$ & 265   & 265 & 15 & $55.82$ & $60.12$ & $738.6$ \\
$z = 50 \rightarrow 10$  & 1{,}591 & 296 & 16 & $25.70$ & $30.75$ & $252.3$ \\
$z = 10 \rightarrow 4$   & 2{,}857 & 153 & 8  & $33.12$ & $35.82$ & $200.1$ \\
$z = 4 \rightarrow 1.5$  & 4{,}430 & 264 & 8  & $63.00$ & $65.82$ & $197.4$ \\
$z = 1.5 \rightarrow 0.5$& 6{,}040 & 384 & 5  & $68.45$ & $70.25$ & $169.2$ \\
$z = 0.5 \rightarrow 0$  & 5{,}011 & 314 & 4  & $58.38$ & $59.90$ & $155.1$ \\
\hline
$z = 200 \rightarrow 0$ & 20{,}194 & 1{,}676 & 56 & $304.48$ & $322.66$ & --- \\
$z = 50 \rightarrow 0$  & 19{,}929 & 1{,}411 & 41 & $248.65$ & $262.54$ & --- \\
\end{tabular}
\end{ruledtabular}
\end{table}

We bracket full-run GPU requirements with two complementary models (\texttt{model\_1024\_full\_history.py}, Fig.~\ref{fig:history} and Table~\ref{tab:fullrun}):
(1) a \textit{Clamped Interval Model} that interpolates measured pure-evolution speedups linearly in scale factor $a$ between $z \approx 2$ ($a=0.346$) and $z \approx 0$ ($a=0.970$) while conservatively capping high-redshift speedup at the $z \approx 2$ anchor ($9.27\times$ for $4\times$ solo), and
(2) a \textit{Step-Physics Regression Model} that evaluates fitted per-step GPU costs $t_{\rm nonPM} = \alpha + \beta (N_{\rm act}/1024^3)$ ($\alpha=1.64$--$4.76\,\mathrm{s}$, $\beta=19.6$--$21.2\,\mathrm{s}$ on $4\times$ solo) and mean PM step times ($20.72$--$22.47\,\mathrm{s}$) across all $20{,}194$ steps, scaled by empirical dynamic-PM step-count ratios ($1.0$ at $z \ge 50$, $910/589$ at $z \approx 2$, $513/757$ at $z \approx 0$). Whole-node electrical energy compares a nominal $700\,\mathrm{W}$ dual-socket 256-core CPU server ($213.1\,\mathrm{kWh}$ pure gravity) against directly measured allocated GPU telemetry ($217.4\,\mathrm{W}$ for $1\times$, $696.5\,\mathrm{W}$ mean solo / $647.1\,\mathrm{W}$ shared across GPUs~\texttt{0..3}) plus $+100$--$120\,\mathrm{W}$ host overhead.

\begin{table}[!tbp]
\caption{Full-history $1024^3$ pure-evolution runtime (hours) and whole-node energy (kWh) projections (\texttt{clamped} / \texttt{step-physics} range).}\label{tab:fullrun}
\scriptsize
\setlength{\tabcolsep}{2pt}
\begin{ruledtabular}
\begin{tabular}{lcccc}
Platform & $z=200{\to}0$ & Speedup & Energy & $z=50{\to}0$ \\
\hline
256-core CPU & $304.5$ ($12.7\,\mathrm{d}$) & $1.00\times$ & $213.1$ & $248.7$ \\
$1\times$ A6000 & $92.9$--$96.1$ & $3.17$--$3.28\times$ & $29.5$--$30.5$ & $76.2$--$91.7$ \\
$4\times$ A6000 solo & $\mathbf{34.6\text{--}35.8}$ & $\mathbf{8.50\text{--}8.80\times}$ & $\mathbf{28.3\text{--}29.2}$ & $\mathbf{28.6\text{--}34.3}$ \\
$4\times$ A6000 shr. & $41.1$--$41.3$ & $7.37$--$7.42\times$ & $31.5$--$31.7$ & $33.6$--$39.7$ \\
\end{tabular}
\end{ruledtabular}
\end{table}

Both independent models agree within $3.5\%$: evolving $1024^3$ particles across the full cosmological history $z = 200 \rightarrow 0$ requires \textbf{$34.6$--$35.8\,\mathrm{h}$ ($1.44$--$1.49$ days)} on four NVLink-paired RTX A6000 GPUs versus $12.69$ node-days of pure gravity on 256 CPU cores ($8.50$--$8.80\times$ speedup), while cutting whole-node energy from $213.1\,\mathrm{kWh}$ down to \textbf{$28.3$--$29.2\,\mathrm{kWh}$ ($\approx 86\%$ savings)}. Even $1\times$ RTX A6000 saves a comparable fraction of electricity ($29.5$--$30.5\,\mathrm{kWh}$) despite taking $3.9$--$4.0$ days. Measured post-processing (\texttt{to-hdf5} averaging $1{,}260\,\mathrm{s/snap}$ and 96-thread CPU FoF at $2{,}986\,\mathrm{s}$ at $z \approx 2$ or $1{,}773\,\mathrm{s}$ at $z=0$) can run asynchronously on host CPUs from binary checkpoints, preserving the $\sim 1.5$-day GPU wall clock for full 56-snapshot campaigns.

\subsection{Cross-Code Runtime, NVML Energy, and Normalized Throughput}

\figwide{fig7_code_energy_comparison.pdf}{Controlled same-node ($2\times$ AMD EPYC 9754 + RTX A6000) four-code comparison on $256^3$ ($a = 0.855 \rightarrow 0.920$) with matched opening angle ($\theta = 0.70$) and timestep tolerance ($\eta_G = 0.025$ / \texttt{PKDGRAV3} $\eta_{\rm PKD} = 0.2236$). (a) Total wall vs.\ pure evolution time. (b) Hardware-measured $200\,\mathrm{ms}$ NVML GPU energy combined with modeled host power ($+120\,\mathrm{W}$ GPU host; utilization-scaled $120$--$800\,\mathrm{W}$ and flat $700\,\mathrm{W}$ CPU reference). (c) Matter power spectrum deviation relative to 256-thread CPU \texttt{SWIFT}, alongside the $1024^3$ full-history energy savings.}{fig:crosscode}

To situate \texttt{gravityCU} against established CPU and GPU cosmological $N$-body solvers under identical hardware and precision settings~\citep{Schneider2016Euclid}, we executed \texttt{CPU SWIFT} ($256$ threads), \texttt{GADGET-4} ($128$ physical EPYC cores, \texttt{PMGRID=128}, TreePM)~\citep{Springel2021GADGET4}, \texttt{PKDGRAV3 3.5.1} (CUDA FMM+Ewald on $1\times$ and $2\times$-MPI RTX A6000)~\citep{Potter2017PKDGRAV3}, and \texttt{gravityCU} ($1\times, 2\times, 4\times$ RTX A6000) on node \texttt{a9} across the $256^3$ window ($a = 0.855084 \rightarrow 0.92$, $\theta = 0.70$, $\epsilon = 0.001953\,h^{-1}\mathrm{Mpc}$, $\eta_G = 0.025$ / \texttt{PKDGRAV3} $\eta_{\rm PKD} = \sqrt{2\eta_G} = 0.2236068$), logging GPU power every $200\,\mathrm{ms}$ via NVML. Because \texttt{a9} disables unprivileged Linux RAPL \texttt{powercap}, CPU/host power is reported transparently as a modeled dual-socket node estimate ($+120\,\mathrm{W}$ host overhead for GPU runs; utilization-scaled $120 + 680 \min(1, u_{\rm CPU}/25600)\,\mathrm{W}$ and a flat $700\,\mathrm{W}$ node reference for CPU runs).

As summarized in Table~\ref{tab:crosscode} and Fig.~\ref{fig:crosscode}, \texttt{gravityCU} sustains \textbf{$18.84\,\mathrm{M}$ active-particle updates/s} on $1\times$ A6000 and \textbf{$49.62\,\mathrm{M}$ updates/s} on $4\times$ A6000 ($12.41\,\mathrm{M/s/GPU}$), finishing the run in $28.26\,\mathrm{s}$ total / $18.99\,\mathrm{s}$ evolution ($28.9\times$ wall / $41.9\times$ evo vs.\ CPU \texttt{SWIFT}) at \textbf{$15.01$--$20.05\,\mathrm{J/M\text{-step}}$} measured GPU energy (\textbf{$3.93$--$5.25\,\mathrm{Wh}$}) and \textbf{$5.27\,\mathrm{Wh}$} minimum whole-node energy on $2\times$ A6000 (\textbf{$25.8\times$} reduction vs.\ $136.16\,\mathrm{Wh}$ CPU \texttt{SWIFT}). All codes satisfy percent-level $P_\delta(k)$ parity ($0.00106\%$--$0.187\%$) and FoF HMF agreement ($91$--$92/92$ at $N\ge 10^4$, $17/17$ at $N\ge 10^5$). Notably, \texttt{PKDGRAV3} at $256^3$ ($16.78\mathrm{M}$ particles) is architecturally bounded by $67$ host-side domain decompositions ($426\,\mathrm{s}$ on $1\mathrm{G}$, $279\,\mathrm{s}$ on $2\mathrm{G}$ MPI) and $245$ CPU tree builds ($522\,\mathrm{s}$ / $146\,\mathrm{s}$) across $244$ adaptive rungs, leaving mean GPU utilization at $3.3$--$5.6\%$ ($\approx 68\,\mathrm{W}$ per A6000, $286$--$374\,\mathrm{J/M\text{-step}}$ GPU) compared to $42.7$--$73.6\%$ GPU utilization in fully VRAM-resident \texttt{gravityCU}.

\begin{table*}[!tbp]
\caption{Controlled same-node (\texttt{a9}) cross-code wall time, normalized active-particle throughput, NVML GPU vs.\ modeled node energy, and $P_\delta(k)$ parity on $256^3$ ($N = 16{,}777{,}216$, $a = 0.855 \rightarrow 0.920$, $\theta = 0.70$; total active updates: $485.66\mathrm{M}$ PKDGRAV3, $942.3$--$942.8\mathrm{M}$ SWIFT/\texttt{gravityCU}, $1034.29\mathrm{M}$ GADGET-4).}\label{tab:crosscode}
\scriptsize
\setlength{\tabcolsep}{2.4pt}
\begin{ruledtabular}
\begin{tabular}{lcccccc}
Code / Config & Wall / Evo (s) & M/s (per GPU) & GPU Wh (W) & Node Wh & J/M-step (GPU / Node) & Max $|\Delta P_\delta|/P_\delta$ \\
\hline
\texttt{PKDGRAV3 3.5.1} (1G) & $2043.2$ / $1926.3^\ddagger$ & $0.252$ ($0.252$) & $38.57$ ($67.9$) & $106.68$ & $285.9$ / $790.8$ & $0.150\%$ ($P_p$: $0.693\%$) \\
\texttt{PKDGRAV3 3.5.1} (2G MPI) & $1333.2$ / $1223.6^\ddagger$ & $0.397$ ($0.198$) & $50.45$ ($136.2$) & $94.89$ & $374.0$ / $703.4$ & $0.150\%$ \\
\texttt{GADGET-4} (128C MPI) & $883.2$ / $857.7$ & $1.206$ (---) & --- & $110.70$ ($171.7^\star$) & --- / $385.3$ ($597.8^\star$) & $0.187\%$ \\
\texttt{CPU SWIFT} (256T base) & $816.0$ / $796.4$ & $1.184$ (---) & --- & $136.16$ ($158.7^\star$) & --- / $519.9$ ($605.9^\star$) & Reference ($92$, $17$) \\
\texttt{gravityCU} 1G (3D Lazy) & $58.09$ / $50.03$ & $18.84$ ($18.84$) & $\mathbf{3.93}$ ($241.3$) & $5.86$ & $\mathbf{15.01}$ / $22.40$ & $0.0387\%$ ($92/92$) \\
\texttt{gravityCU} 2G (3D Lazy) & $33.11$ / $28.66$ & $32.88$ ($16.44$) & $4.17$ ($448.1$) & $\mathbf{5.27}$ & $15.93$ / $\mathbf{20.15}$ & $0.0387\%$ ($92/92$) \\
\texttt{gravityCU} 4G (3D Eager) & $30.21$ / $21.07$ & $44.74$ ($11.19$) & $5.38$ ($635.2$) & $6.39$ & $20.54$ / $24.39$ & $\mathbf{0.00106\%}$ ($92/92$) \\
\texttt{gravityCU} 4G (3D Lazy) & $\mathbf{28.26}$ / $\mathbf{18.99}$ & $\mathbf{49.62}$ ($12.41$) & $5.25$ ($662.7$) & $6.19$ & $20.05$ / $23.64$ & $0.0387\%$ ($92/92$) \\
\end{tabular}
\end{ruledtabular}
\begin{flushleft}
\footnotesize $^\ddagger$Excludes tipsy/HDF5 read and write I/O ($116.9\,\mathrm{s}$ on $1\mathrm{G}$, $109.6\,\mathrm{s}$ on $2\mathrm{G}$ MPI); throughput uses non-I/O evolution time. $^\star$Parentheses show flat $700\,\mathrm{W}$ dual-socket CPU node reference alongside utilization-scaled EPYC 9754 model.
\end{flushleft}
\end{table*}

\section{Conclusions and Code Availability}\label{sec:conc}

We have demonstrated that keeping $1024^3$ self-gravitating cosmological particles 100\% resident in VRAM across four $48\,\mathrm{GiB}$ GPUs completely eliminates host particle swapping while achieving $7.26\times$--$9.27\times$ acceleration over 256-core CPU \texttt{SWIFT} ($0.0045\%$--$0.010\%$ $P_\delta$, $0.011\%$--$0.041\%$ aligned $P_\theta$, $\pm 0.32\%$ FoF HMF) and projecting full $z = 200 \rightarrow 0$ campaigns in $\approx 1.45\,\mathrm{days}$ and $28.3$--$29.5\,\mathrm{kWh}$. Extending the solver with lazy Hilbert octree refits, split zero-spill kernels, distributed cuFFT PM, $12\times 12\times 12$ 3D domain decomposition, and Full-GPU FoF/Catalog/SO boosts $256^3$ execution to $17.53\,\mathrm{s}$ ($45.4\times$ evo, $49.62\,\mathrm{M}$ updates/s) and $5.27$--$6.19\,\mathrm{Wh}$ whole-node energy ($22.0$--$25.8\times$ savings vs.\ CPU \texttt{SWIFT} and $18.0\times$ vs.\ \texttt{PKDGRAV3}), while passing Zel'dovich plane-wave, exact Ewald $\theta$-convergence, Plummer virial, and Layzer--Irvine physics gates. Resolving EPYC PCIe \texttt{NODE} peer-access deadlocks through native-atomic NVLink gating (\texttt{GRAV\_P2P\_MODE=nvlink}) provides a clean single-node foundation for multi-node MPI + GPU-Direct RDMA scaling. All CUDA kernels, evaluation scripts, verification JSON/CSV datasets, and figures are open-source at \url{https://github.com/HainaHuang729/grav-gpu}.

\begin{acknowledgments}
GPU benchmarks and end-to-end validations were executed on the Department of Physics GPU cluster (\texttt{a9}) at The Chinese University of Hong Kong, and reference archival snapshots were produced on the CUHK Central High-Performance Computing Cluster (\texttt{PL\_25\_1024}).
\end{acknowledgments}

\bibliography{references}

\end{document}
